\section{Conclusion}

This work presents {TemporalBench}, a benchmark for systematically evaluating agent intelligence on time-series tasks that require temporal understanding, prediction, and contextual reasoning. By combining real-world time-series data with event-centric task construction and multi-tier evaluation (T1--T4), the benchmark moves beyond traditional forecasting metrics and exposes how agents reason, fail, and generalize under realistic temporal decision settings.

Our empirical results reveal several consistent patterns. First, strong numerical forecasting performance does not translate into reliable temporal reasoning or correct decision-making in discrete task settings. Across agents and domains, reasoning-intensive tasks---especially those involving contextual interpretation or event-aware prediction---remain challenging. Second, many failures stem from execution and controllability issues rather than an absence of temporal signal. Errors related to forecast horizon length, output formatting, and numerical instability dominate across agent frameworks, even when all systems share the same backbone model. Third, increasing historical input length does not guarantee improved performance. Instead, length sensitivity is highly dataset- and task-dependent, with longer contexts often introducing instability rather than consistent gains.

Our analyses further suggest that current agent designs exhibit fragmented strengths. General-purpose agents show limited sensitivity to additional temporal context, single-LLM baselines display clear trade-offs between contextual richness and controllability, and time-series--specific agents can exploit longer histories in some cases but remain brittle in others. Lightweight representation augmentations, such as feature extraction and visualization, can alleviate certain failure modes, but are insufficient to address the broader challenges of temporal reasoning and robust structured generation.

Taken together, these findings indicate that progress in agentic time-series modeling will require more than stronger language models or longer prompts. Future agents must incorporate explicit mechanisms for temporal abstraction, history selection, and output validation, treating execution constraints as first-class objectives alongside reasoning accuracy. From a benchmarking perspective, {TemporalBench} highlights the importance of evaluating agents across multiple task types, domains, and input regimes, and provides a foundation for diagnosing where agent intelligence succeeds---and where it still fundamentally falls short---in temporal decision-making.



